% ══════════════════════════════════════════════════════════════════════
% Appendix AI: Development Methodology
% ══════════════════════════════════════════════════════════════════════
\section{Development Methodology}\label{app:dev-methodology}

This appendix documents the progressive development methodology used to
build the SAR drone system, from initial simulation through to field
integration.  It covers the design philosophy, a 10-phase timeline with
outcomes, and lessons that future teams can apply.

% ──────────────────────────────────────────────────────────────────────
\subsection{Development Philosophy}\label{sec:dev-philosophy}

Three principles governed development from the outset:

\begin{enumerate}[leftmargin=1.5em]
  \item \textbf{Simulation-first.}  Every algorithmic component---state
        machine, search pattern generator, detection pipeline, GPS
        estimation---was prototyped and end-to-end tested in SITL
        (Software-In-The-Loop) on a laptop before any hardware was
        touched.  SITL provides repeatable conditions, instant restart,
        and zero crash risk.

  \item \textbf{Same code everywhere.}  A single codebase runs on the
        Windows development laptop, WSL/Linux, and the Raspberry~Pi~5.
        Platform differences are confined to \texttt{config.py}
        (auto-detection of serial ports, camera backends, display
        availability) and \texttt{vision.py} (dual backend:
        Ultralytics on laptop, TFLite/NCNN on Pi).  No
        \texttt{if platform ==} guards appear in mission logic.

  \item \textbf{Progressive trust.}  Each test step isolates exactly
        one new variable.  Bench tests precede flight tests; passive
        observation precedes autonomous commands; logging-only modes
        precede action-taking modes.  A failure at step~$N$ is
        attributable to the single change introduced at step~$N$.
\end{enumerate}

% ──────────────────────────────────────────────────────────────────────
\subsection{Ten-Phase Development Timeline}\label{sec:dev-timeline}

\tabref{tab:dev-phases} summarises the ten development phases, their
scope, test outcomes, and calendar dates.

\begin{table}[H]
\centering
\compactTable
\caption{Ten-phase development history with outcomes.  The majority of integration defects were discovered in phases~5--7 (bench and field testing), validating the progressive approach.}
\label{tab:dev-phases}
\begin{tabularx}{\linewidth}{@{}c l X l@{}}
\toprule
\textbf{\#} & \textbf{Phase} & \textbf{Scope and Outcomes} & \textbf{Dates} \\
\midrule
1 & Simulation on Laptop &
    State machine (initial 11~states, later expanded to 20), lawnmower search pattern,
    simulated drone camera from \texttt{map.jpg}, Ultralytics YOLOv8n
    detection, GPS estimation, operator Y/N confirmation, SITL
    integration.  Full mission completed end-to-end in simulation. &
    Weeks 12--13 \\
\addlinespace
2 & Multi-Platform CV &
    Dual-backend \texttt{vision.py}: Ultralytics (laptop) / TFLite
    (Pi).  Docker Pi-environment test passed.  TFLite coordinate bug
    fixed (pixel vs.\ normalised).  \texttt{config.py} auto-detects
    connection type. &
    Week 13 \\
\addlinespace
3 & Test Infrastructure &
    \texttt{preflight.py}, 13~test scripts covering camera, Cube
    heartbeat, GPS, AI detection, FOV calibration, resolution trade-off.
    Organised into \texttt{tests/hardware/}, \texttt{tests/flight/},
    \texttt{tests/calibration/}, \texttt{tests/diagnostics/}. &
    Week 14 \\
\addlinespace
4 & Pi + Camera &
    Raspberry~Pi~5 with IMX296 global shutter.  Camera verified at
    $1456\times1088$.  BGR colour issue discovered and fixed (LL-02).
    TFLite inference: \SI{206.5}{\milli\second} average, 4.8~FPS,
    0.966~confidence (50-run benchmark). &
    Week 16 \\
\addlinespace
5 & Pi + Cube Wired &
    Cube Orange connected via UART at \SI{921600}{baud}.  Mavproxy
    UDP bridge resolves Python~3.13 serial bug (LL-07).  Mission
    Planner connects via TCP bridge. &
    Week 16 \\
\addlinespace
6 & Interactive Simulator &
    \texttt{simple\_simulator.py}: keyboard flight, visual servo,
    GPS estimation, spatial clustering, Kalman filter,
    7.5\,m offset landing, multi-target support, data recording. &
    Weeks 14--15 \\
\addlinespace
7 & Web Ground Station &
    \texttt{pi\_flight.py}: MJPEG stream, 2D GPS grid, browser
    dashboard with detection clusters and command buttons.
    Headless over SSH (PuTTY).  Tested in SIMULATION on laptop. &
    Week 17 \\
\addlinespace
8 & Field Integration &
    FOV calibrated (\SI{5.46}{\milli\metre} focal length at bench;
    \SI{54.4}{\degree} HFOV from DJI video at altitude).  Lens
    calibration (RMS~0.399).  Pi bench testing with real Cube.
    Weather cancelled outdoor flight; adapted to bench-only protocol. &
    Weeks 19--20 \\
\addlinespace
9 & Model Retraining &
    \texttt{generate\_dataset\_v2.py}: 366~images (300~synthetic +
    16~real-labelled + 50~negatives) at $1456\times1088$.
    Trained YOLOv8n on Colab: \texttt{imgsz=1088}, 150~epochs,
    mAP50\,=\,0.995.  NCNN backend added (72\,ms, 4.5$\times$
    faster than TFLite). &
    Week 18 \\
\addlinespace
10 & Safety \& Polish &
    RC override guard (LL-01).  Five-layer SSSI geofence (speed field, repulsion, hard cutoff,
    waypoint filter, firmware backup).  SmartEstimator greedy-cluster GPS lock.  Threaded
    inference (stream at 30\,fps, detect in background).
    127~unit tests, 33~automated integration tests. &
    Weeks 19--20 \\
\bottomrule
\end{tabularx}
\end{table}

% ──────────────────────────────────────────────────────────────────────
\subsection{Simulation Stage (Weeks 12--15)}\label{sec:sim-stage}

All mission logic was developed and validated in SITL before any
hardware was procured.  The simulation environment consisted of:

\begin{itemize}[leftmargin=1.5em]
  \item \textbf{ArduPilot SITL} via Mission Planner, providing
        realistic GPS, attitude, and flight dynamics.
  \item \textbf{Simulated camera} (\texttt{simulation.py}): a
        satellite image (\texttt{map.jpg}) with a virtual drone viewport
        that moves according to SITL telemetry.
  \item \textbf{State machine} (\texttt{main.py}): 20~states from
        \texttt{INIT} through \texttt{DONE}, with 32~transitions
        handling detection, centering, descent, verification, and
        landing.
  \item \textbf{Search pattern} (\texttt{planning.py}): lawnmower
        generator for arbitrary convex polygons with configurable
        strip spacing, edge margins, and start-corner optimisation.
  \item \textbf{Interactive simulator}
        (\texttt{simple\_simulator.py}, 2508~lines): keyboard-driven
        flight with GPS estimation, inverse-variance weighting, Kalman
        filtering, spatial clustering, and 7.5\,m offset landing.
\end{itemize}

By the end of the simulation stage the full mission---takeoff, search,
detect, centre, descend, verify, land---ran end-to-end on the laptop
with zero hardware dependencies.

% ──────────────────────────────────────────────────────────────────────
\subsection{Pi Deployment (Weeks 16--18)}\label{sec:pi-deploy}

Hardware integration proceeded component-by-component:

\begin{enumerate}[leftmargin=1.5em]
  \item \textbf{Camera.}  IMX296 global shutter at
        $1456\times1088$ native resolution.  A persistent blue tint
        was traced to the sensor outputting BGR despite picamera2's
        \texttt{RGB888} label; removing \texttt{cv2.cvtColor} fixed it
        (LL-02).

  \item \textbf{AI inference.}  TFLite on Pi~5 CPU (XNNPACK):
        \SI{206.5}{\milli\second} per frame, 4.8~FPS, 50/50
        detections at 0.966~confidence (\texttt{tests/hardware/%
        benchmark.py}).  NCNN backend later added:
        \SI{72}{\milli\second}, 9~FPS full pipeline.

  \item \textbf{Cube connection.}  Python~3.13 breaks
        \texttt{pyserial} reads; mavproxy UDP bridge
        (\texttt{udpout:127.0.0.1:14550}) resolves this with zero
        data loss (LL-07).  Mission Planner connects via
        \texttt{tcpin:0.0.0.0:5762}.

  \item \textbf{Model retraining.}  \texttt{generate\_dataset\_v2.py}
        produced 366~images at native resolution.
        \texttt{tools/label\_tool.py --full} labelled 16~real frames
        from DJI video.  Colab training yielded mAP50\,=\,0.995;
        the model is a drop-in replacement (same input/output shapes).
\end{enumerate}

% ──────────────────────────────────────────────────────────────────────
\subsection{Field Integration (Weeks 19--20)}\label{sec:field-int}

A field day at the designated site (2026-03-11) assembled the full
Pi + Cube + camera stack.  Weather cancelled the planned flight;
the team adapted by executing every possible bench and calibration
activity:

\begin{itemize}[leftmargin=1.5em]
  \item \textbf{FOV calibration:} \SI{92}{\centi\metre} visible at
        \SI{1}{\metre} height $\Rightarrow$
        $f = \SI{5.46}{\milli\metre}$ (updated from the
        initial 7.0\,mm estimate).
  \item \textbf{Lens calibration:} checkerboard method (14$\times$9,
        28\,mm squares), RMS\,=\,0.399, undistortion cost
        $\sim$\SI{1.5}{\milli\second} per frame.
  \item \textbf{Connectivity:} three-terminal PuTTY workflow verified
        (mavproxy, diagnostics, flight script).  Port conflicts and
        camera-ownership issues resolved.
  \item \textbf{Bug discovery:} corrupt
        \texttt{calibration\_data.npz} (0~bytes) was mangling frames
        via \texttt{cv2.remap}; deleting it restored detection
        immediately (LL-04).
  \item \textbf{Subsequent video analysis} of DJI flight footage
        revealed a 100--200\,ms GPS timing lag causing
        CEP50\,=\,\SI{2.3}{\metre} scatter in target estimates
        (LL-05).
\end{itemize}

% ──────────────────────────────────────────────────────────────────────
\subsection{Ambition Levels}\label{sec:ambition-levels}

The project defined three deliverable tiers to manage scope risk (Table~\ref{tab:ambition-levels}):

\begin{table}[H]
\centering
\compactTable
\caption{Three product levels with autonomy and demonstration scope.  The Minimum Viable Deliverable (Level~1) was achieved; Levels~2 and~3 remain contingent on successful flight testing.}
\label{tab:ambition-levels}
\begin{tabularx}{\linewidth}{@{}c l l X@{}}
\toprule
\textbf{Level} & \textbf{Name} & \textbf{Autonomy} & \textbf{Demonstration Scope} \\
\midrule
L1 & Manual MVP & None &
  Pilot flies RC manually.  Pi runs camera + AI passively
  (zero commands).  Detection overlay on video stream, buzzer alerts,
  geotagged images saved.  Pilot lands manually. \\
\addlinespace
L2 & Semi-Autonomous & Partial &
  Drone arms, takes off, flies lawnmower search avoiding SSSI.
  CV detects targets; pilot classifies (Y/I/X).  On confirm: centre,
  descend, offset land at 7.5\,m, payload release.  RC override
  always available. \\
\addlinespace
L3 & Full Autonomous & Full &
  As L2 but confidence accumulation replaces pilot confirmation.
  Multi-pass verification at decreasing altitude.  Drone auto-decides. \\
\bottomrule
\end{tabularx}
\end{table}

This tiered approach ensured that even if weather or hardware issues
prevented full autonomous flight, L1 would still produce a
demonstrable, useful system and collect the data needed for L2
development.

% ──────────────────────────────────────────────────────────────────────
\subsection{Progressive Testing Ladder}\label{sec:test-ladder}

Flight testing follows the strict five-step ladder shown in Table~\ref{tab:test-ladder}.  Each step isolates
one new risk factor; no step may be skipped.

\begin{table}[H]
\centering
\compactTable
\caption{Progressive flight testing ladder with scripts and risk isolation.  Each step introduces exactly one new risk factor; failures must be resolved before advancing.}
\label{tab:test-ladder}
\begin{tabularx}{\linewidth}{@{}c X l l@{}}
\toprule
\textbf{Step} & \textbf{Description} & \textbf{Script} & \textbf{New Risk} \\
\midrule
0 & Bench commands: arm, mode change, servo PWM.
    No propellers. &
  \texttt{0a--0f\_*.py} & Wiring \\
\addlinespace
1 & Manual RC flight.  Pi runs camera + AI, logs detections,
    buzzer beeps.  \textbf{Zero MAVLink commands.} &
  \texttt{1\_passive\_flight.py} & Outdoor CV \\
\addlinespace
2 & Code-controlled waypoint flight (3--4 GPS points).
    No camera, no detection, no decisions. &
  \texttt{2\_waypoint\_test.py} & MAVLink flight \\
\addlinespace
3 & AUTO waypoints + AI detection.  On detection, switch
    to GUIDED hover.  No descent or landing. &
  \texttt{3\_auto\_detect.py} & CV + flight \\
\addlinespace
4 & Full autonomous: search pattern, detect, centre,
    descend, verify (Y/N), offset land. &
  \texttt{4\_detect\_and\_center.py} & Full mission \\
\bottomrule
\end{tabularx}
\end{table}

Six bench-level tests (step~0) additionally cover planning validation
(\texttt{0d\_planning\_test.py}, 6~unit tests), geofence sign
conventions (\texttt{0e\_geofence\_test.py}, 7~unit tests), and
payload servo actuation (\texttt{0f\_servo\_test.py}).

% ──────────────────────────────────────────────────────────────────────
\subsection{Achievement Summary}\label{sec:achievement-summary}

An objective scoring matrix (\tabref{tab:sim-vs-real}) was constructed
for each requirement across both simulation and real hardware.

\begin{table}[H]
\centering
\compactTable
\caption{Simulation vs.\ real hardware readiness (0--10 scale).  Simulation scores average 7.7/10; real hardware averages 2.1/10, reflecting the absence of outdoor flight testing.}
\label{tab:sim-vs-real}
\begin{tabular}{@{}l c c@{}}
\toprule
\textbf{Category} & \textbf{Simulation} & \textbf{Real Hardware} \\
\midrule
All requirements (R01--R10) & 7.7 & 2.1 \\
Flight-critical (R01, R02, R03, R04, R09) & 8.4 & 1.8 \\
Mission-critical (R05, R06, R07) & 7.0 & 1.0 \\
Infrastructure (R08, R10) & 7.0 & 4.5 \\
\bottomrule
\end{tabular}
\end{table}

\textbf{Simulation (7.7/10):}  The state machine logic, search pattern
generation, CV detection pipeline, GPS estimation mathematics, and
safety systems (RC override, geofence, RTL) all work correctly in
simulation.  The full mission has been completed end-to-end multiple
times in SITL.

\textbf{Real hardware (2.1/10):}  Individual components are verified
(camera: working; TFLite: 4.8~FPS at 0.966~confidence; Cube
connection: stable via mavproxy), but zero minutes of real flight
time have been logged.  The 5.6-point gap between simulation and
real hardware is the project's largest risk.

\textbf{What simulation proved:}  Algorithm correctness, state
transitions, search coverage, detection-to-decision loop, safety
logic.

\textbf{What simulation cannot prove:}  Whether the drone physically
flies; whether CV detects from 15--30\,m altitude outdoors; whether
GPS accuracy is sufficient for offset landing; whether the video
stream is usable at operational WiFi range; whether the geofence
actually prevents SSSI entry in wind.

% ──────────────────────────────────────────────────────────────────────
\subsection{Lessons for Future Teams}\label{sec:lessons-future}

Eight lessons emerged from the development process, each documented
with root cause, fix, and checklist (see Appendix~AH for the full
defect register):

\begin{enumerate}[leftmargin=1.5em]
  \item \textbf{Start hardware testing on day one.}  The simulation
        stage produced mature, well-tested code, but the 5.6-point
        sim-to-real gap shows that software maturity does not transfer
        to hardware confidence.  Even basic tasks---outdoor GPS fix,
        one manual hover---would have closed multiple requirements.

  \item \textbf{Validate the environment before the field.}  Docker
        testing caught missing packages early, but the corrupt
        \texttt{calibration\_data.npz} and BGR colour issue were
        discovered only on-site.  A pre-field smoke test with the exact
        Pi image would have saved hours.

  \item \textbf{Plan for weather cancellation.}  The March~11 field
        day was the sole outdoor window; weather cancelled it.  Having a
        bench-only fallback protocol (FOV calibration, lens calibration,
        connectivity verification) ensured the day was still productive.

  \item \textbf{Decouple camera ownership.}  Only one process can open
        picamera2.  This forced a clean separation: one script owns the
        camera and streams to all consumers via MJPEG.  Design for
        single-owner resources from the start.

  \item \textbf{Never trust format labels.}  The IMX296 sensor
        outputs BGR despite \texttt{RGB888} labelling.  YOLOv8 TFLite
        outputs pixel coordinates despite documentation suggesting
        normalised values.  Test empirically; label after measurement.

  \item \textbf{Build zero-command observation scripts.}  The passive
        watch approach (\texttt{passive\_watch.py}, zero MAVLink
        commands) was the safest way to validate CV in the field.  It
        could run during any flight without risk, collecting training
        data and detection statistics simultaneously.

  \item \textbf{Use progressive trust, never skip steps.}  The
        numbered test ladder (0a--4) isolates exactly one new variable
        per step.  When Step~$N$ fails, the cause is the single change
        introduced at Step~$N$, not accumulated uncertainty from
        earlier steps.

  \item \textbf{Maintain tiered deliverables.}  L1~(Manual MVP),
        L2~(Semi-Autonomous), L3~(Full Autonomous) ensured that even
        if only L1 was achievable, it would be a complete, demonstrable
        product rather than a half-finished L2.
\end{enumerate}
